\begin{tikzpicture}[scale=1.0, font=\small]
  % ---------- small-signal equivalent ----------
  \begin{scope}
    \node[note, anchor=south] at (3.2,3.4){共射级的小信号等效电路};
    \draw (0,2.4) node[left]{$v_{in}$} -- (1.1,2.4) node[circ]{};
    \draw (1.1,2.4) to[R=$r_\pi$] (1.1,0);
    \draw (1.1,0) -- (5.9,0);
    \node[ground] at (3.5,0){};

    \draw (1.1,2.4) -- (3.1,2.4);
    \draw (3.1,2.4) node[circ]{} to[cI, l_=$g_mv_\pi$, invert] (3.1,0);
    \draw (4.3,2.4) node[circ]{} to[R=$r_o$] (4.3,0);
    \draw (5.5,2.4) node[circ]{} to[R=$R_C$] (5.5,0);
    \draw (3.1,2.4) -- (6.6,2.4) node[right]{$v_{out}$};

    \draw[helper] (0.55,2.55) -- (0.55,0.7);
    \node[note, anchor=east, font=\footnotesize] at (0.5,1.6){$v_\pi$};
  \end{scope}

  \node[note, anchor=north, align=left] at (0.2,-0.9)
    {输出节点的 KCL：$\dfrac{v_{out}}{r_o\parallel R_C}=-g_mv_\pi$，而 $v_\pi=v_{in}$};

  \node[note, anchor=north west, align=left, text=ecblue] at (7.8,3.3)
    {$A_v=-g_m\,(R_C\parallel r_o)$\\[6pt]
     $R_{in}=r_\pi=\dfrac{\beta}{g_m}$\\[6pt]
     $R_{out}=R_C\parallel r_o$};

  \node[note, anchor=north west, align=left] at (7.8,0.6)
    {把 $g_m=I_C/V_T$ 代进去：\\[3pt]
     $|A_v|=\dfrac{I_CR_C}{V_T}=\dfrac{V_{R_C}}{V_T}$\\[5pt]
     增益 = 「$R_C$ 上的直流压降 / \SI{26}{mV}」——\\
     所以增益上限被电源电压卡死，\\
     这是共射级最重要的一条设计直觉};
\end{tikzpicture}
